% Ad Familiares 16.2 --- headnote
% ad-familiares-16-02, 5 November 50 BC, Alyzia

\textit{Cicero to Tiro, written from Alyzia on the
Nones of November 50 BC --- 5 November --- two days
after the previous letter from the same voyage. Cicero
has come ashore at the small Acarnanian port a hundred
and twenty stadia south of Leucas; he had expected to
find Tiro himself, or at least a letter through the
slave Marion, waiting at the next port north, and the
silence has put him in the state of mind he refuses to
describe.}

\textit{The note is the shortest and most uneasy of the
cluster --- a single section under twenty lines of Latin,
the salutation now stripped down to \textit{Tullius
Tironi suo s.} (the family chorus of 16.1 dropped). The
opening is the rhetorical gesture of a man writing
against his own intention: \textit{non queo ad te nec
libet scribere quo animo sim adfectus} --- I cannot,
and do not wish to, tell you what state I am in --- and
the close folds the formal love-formula in on itself,
asking Tiro to measure his care for his health either by
his love for Cicero or by his knowledge of Cicero's love
for him, whichever weighs more.}
